\chapter{Complex Numbers and Roots of Unity}\label{ch:complex}

\begin{center}\itshape
Why is $\cos\frac{2\pi}5=\frac{\sqrt5-1}4$, and why does $\sin\frac\pi7\,\sin\frac{2\pi}7\cdots\sin\frac{6\pi}7=\frac7{64}$?\\
Both answers live on the unit circle.
\end{center}

Both numbers are real. The first is the cosine of one fifth of a full turn. The second is a product of six sines. Nothing in either formula mentions a circle, and both formulas are bookkeeping for points on the unit circle. The point $e^{2i\pi/5}$ is a vertex of a regular pentagon. The point $e^{2i\pi/7}$ is a vertex of a regular heptagon. Once multiplication means ``stretch and turn'', the two identities are computations with those vertices.

The cancellation is exact. Add the vertices of a regular polygon as arrows from the origin, and the vector sum is $0$. A random handful of arrows of the same length would wander. These close.

Complex numbers were forced on mathematicians by equations whose answers were real. In $1545$ Gerolamo Cardano published the \emph{Ars Magna}. It contains the solution of the cubic, and a quarrel. In $1539$ he had promised Tartaglia not to publish Tartaglia's method. He published a method in $1545$, after learning that Scipione del Ferro had a cubic formula earlier. In Chapter~37 he asks for two numbers with sum $10$ and product $40$. The formal answer is $5+\sqrt{-15}$ and $5-\sqrt{-15}$. Their sum is $10$. Their product is
\[
\bigl(5+\sqrt{-15}\bigr)\bigl(5-\sqrt{-15}\bigr)=25-(-15)=40.
\]
He checks this with the phrase \emph{dimissis incruciationibus}. Witmer's translation is ``putting aside the mental tortures involved''; the same Latin has been read as the cancellation of the cross terms. Cardano then calls the calculation as subtle as it is useless, and moves on.

Rafael Bombelli did not move on. He was an engineer, and he worked on draining the marshes of the Val di Chiana. Cardano's formula for $x^3=px+q$ reads
\[
x=\sqrt[3]{\frac q2+\sqrt{\frac{q^2}4-\frac{p^3}{27}}}+\sqrt[3]{\frac q2-\sqrt{\frac{q^2}4-\frac{p^3}{27}}}.
\]
For $x^3=15x+4$ the root $x=4$ is plain to see, and the formula produces $\sqrt[3]{2+\sqrt{-121}}+\sqrt[3]{2-\sqrt{-121}}$. In his \emph{Algebra} of $1572$ Bombelli wrote that this sort of root looked ``pi\`u tosto sofistica che reale'', more sophistic than real. He calculated with it anyway. He laid down rules for ``pi\`u di meno'' and ``meno di meno'', plus of minus and minus of minus. He found $(2+\sqrt{-1})^3=2+\sqrt{-121}$, and the formula gave
\[
\bigl(2+\sqrt{-1}\bigr)+\bigl(2-\sqrt{-1}\bigr)=4.
\]
The impossible quantities cancelled, and the real root came out. He did not yet have the other two real roots. They are $-2\pm\sqrt3$, and Exercise~\ref{ex:cardano} recovers all three, once a number is allowed three cube roots rather than one.

The suspicion lasted two more centuries. On $5$ May $1777$ Euler presented a memoir in which $\sqrt{-1}$ is written $i$, with $i^2=-1$. The memoir was printed in $1794$. On $10$ March $1797$ Caspar Wessel, a Norwegian surveyor, presented the Royal Danish Academy with a paper that draws these numbers as arrows in the plane: multiplying is stretching and turning. The paper was printed in the Academy's memoirs for $1799$, and it sat unread. It was noticed again in $1895$, and a French translation appeared in $1897$. The picture is usually named after Jean-Robert Argand, a Paris accountant, who circulated the same idea in $1806$. A journal version appeared in Gergonne's \emph{Annales} in $1813$. The polynomials chapter returns to his argument that every nonconstant polynomial has a root.

Gauss had already used the plane in his $1799$ dissertation, which argued for that same existence theorem. In $1831$ he gave the numbers their name, complex numbers. Long before the name, he had the most famous consequence. On $30$ March $1796$, aged eighteen, he wrote in his diary that the regular $17$-gon can be constructed with ruler and compass. The regular polygons known from antiquity were a short list. The $17$-gon was not on it. The common secret, in Bombelli's cube and in Gauss's polygon, is that a turn has become a multiplication, and a regular polygon has become the solution set of $z^n=1$.

\portraits{cardano}{Gerolamo Cardano (1501--1576). Woodcut, eighteenth century. Public domain, via Wikimedia Commons.}{euler}{Leonhard Euler (1707--1783). Portrait by Jakob Emanuel Handmann, 1753, Kunstmuseum Basel. Public domain, via Wikimedia Commons.}

\section{Complex numbers from scratch}

A complex number has to be something we can add and multiply, with $i^2=-1$, and with a reciprocal for every number except $0$. The cheapest way to arrange that is to stop treating $i$ as a mystery and to treat $x+iy$ as the pair $(x,y)$.

\begin{definitionx}
A complex number is a pair of real numbers, written $z=x+iy$. The real part is $x=\operatorname{Re}z$ and the imaginary part is $y=\operatorname{Im}z$. Addition is coordinate by coordinate,
\[
(x+iy)+(x'+iy')=(x+x')+i(y+y').
\]
Multiplication uses the single extra rule $i^2=-1$:
\[
(x+iy)(x'+iy')=(xx'-yy')+i(xy'+yx').
\]
The \emph{conjugate} of $z$ is $\bar z=x-iy$. The \emph{modulus} is $|z|=\sqrt{x^2+y^2}$.
\end{definitionx}

The modulus is the distance from $(x,y)$ to the origin, so $|z|\ge0$, and $|z|=0$ only for $z=0$. Also $|z|=\sqrt{z\bar z}$, because $z\bar z=(x+iy)(x-iy)=x^2+y^2$.

Addition is the addition of vectors in the plane. It is commutative and associative because addition of real numbers is. The same is true of multiplication: the only new ingredient is $i^2=-1$, and that rule is used the same way on either side of a product. Multiplication distributes over addition for the same reason it does for pairs of real numbers. The zero is $0+0i$, and the one is $1+0i$.

The real numbers sit inside $\C$ as the numbers with imaginary part $0$. For those, the product above is the ordinary product. We write $\C^*=\C\setminus\{0\}$.

Every $z\ne0$ has a reciprocal. The computation is
\[
(x+iy)\cdot\frac{x-iy}{x^2+y^2}=\frac{x^2+y^2}{x^2+y^2}=1,
\]
so
\[
\frac1z=\frac{\bar z}{|z|^2}=\frac{x-iy}{x^2+y^2}.
\]
Thus $\C$ is a field.

Conjugation respects the operations. From the formulas, $\overline{z+w}=\bar z+\bar w$ and $\overline{zw}=\bar z\,\bar w$. A number equals its conjugate exactly when its imaginary part is $0$, so $z$ is real if and only if $\bar z=z$. Taking moduli in the product rule gives the identity that makes lengths easy.

\begin{propositionx}[Brahmagupta's identity]
For $z,w\in\C$, $|zw|=|z|\,|w|$. In coordinates, if $z=a+ib$ and $w=c+id$, then
\[
(a^2+b^2)(c^2+d^2)=(ac-bd)^2+(ad+bc)^2.
\]
\end{propositionx}
\begin{proof}
$|zw|^2=zw\,\overline{zw}=zw\,\bar z\,\bar w=|z|^2|w|^2$. Both moduli are nonnegative, so $|zw|=|z|\,|w|$. The coordinate form is the same identity, because $|z|^2=a^2+b^2$, $|w|^2=c^2+d^2$, and $zw=(ac-bd)+i(ad+bc)$.
\end{proof}

Brahmagupta recorded the coordinate form in $628$, as a rule for composing sums of two squares. Check it on $1+2i$ and $2+3i$. The left side is $(1+4)(4+9)=65$. The right side is $(2-6)^2+(3+4)^2=16+49=65$. A product of two sums of two squares is again a sum of two squares. The inequalities chapter uses that fact under this name.

The pair $(x,y)$ is a point of the plane. The number $|z-w|$ is the distance between $z$ and $w$. The inequality that governs estimates is the triangle inequality: a straight path is at most as long as a broken one.

\begin{figure}[ht]
\centering
\fig{ch02-triangle}
\caption{The arrow $z+w$ is one side of a triangle. The other two sides are a copy of $z$ and a copy of $w$. The straight side is at most the sum of the other two, and it is equal only when $z$ and $w$ point along the same ray from the origin.}
\label{fig:triangle}
\end{figure}

\begin{propositionx}[triangle inequality]
For $z,w\in\C$, $|z+w|\le|z|+|w|$, with equality if and only if one of $z,w$ is a nonnegative real multiple of the other. Consequently $\bigl||z|-|w|\bigr|\le|z-w|$.
\end{propositionx}
\begin{proof}
Expand the square of the left side:
\[
|z+w|^2=(z+w)(\bar z+\bar w)=|z|^2+|w|^2+z\bar w+\overline{z\bar w}=|z|^2+|w|^2+2\operatorname{Re}(z\bar w).
\]
The real part of a complex number is at most its modulus, because if $u=s+it$ then $s\le\sqrt{s^2+t^2}$. Apply that to $u=z\bar w$:
\[
\operatorname{Re}(z\bar w)\le|z\bar w|=|z|\,|w|,
\]
where the last step is Brahmagupta's identity. Therefore
\[
|z+w|^2\le|z|^2+|w|^2+2|z|\,|w|=(|z|+|w|)^2.
\]
Both sides are nonnegative, so $|z+w|\le|z|+|w|$.

Equality holds exactly when $\operatorname{Re}(z\bar w)=|z\bar w|$, which means that $z\bar w$ is real and nonnegative. If $w=0$, then $w=0\cdot z$. If $w\ne0$, set $\lambda=z\bar w/|w|^2$. Then $\lambda\ge0$ and $z=\lambda w$. Conversely, if $z=\lambda w$ or $w=\lambda z$ with $\lambda\ge0$, the two arrows point along the same ray, and $|z+w|$ is $(\lambda+1)$ times the modulus of the other factor, which is $|z|+|w|$.

For the consequence, $|z|=|(z-w)+w|\le|z-w|+|w|$, so $|z|-|w|\le|z-w|$. The same bound with $z$ and $w$ swapped gives $|w|-|z|\le|z-w|$.
\end{proof}

Here is a check. Let $z=3+4i$ and $w=6+8i$. Then $w=2z$ with $2\ge0$, and $|z+w|=|9+12i|=15$, while $|z|+|w|=5+10=15$. Replace $w$ by $-z$. Then $|z+w|=0$ and $|z|+|w|=10$, and the inequality is strict, because $-1$ is not a nonnegative multiplier.

\begin{examplex}[square roots by algebra]
Does $i$ have a square root? We want $(x+iy)^2=i$, so $x^2-y^2=0$ and $2xy=1$. The first equation says $y=\pm x$. The second then forces the same sign, so $y=x$ and $2x^2=1$. Thus $x=\pm1/\sqrt2$ and
\[
\sqrt i=\pm\frac{1+i}{\sqrt2}.
\]
Square the right side if you want the check: $\bigl((1+i)/\sqrt2\bigr)^2=(1+2i-1)/2=i$.

The same two equations find the square roots of any $a+ib$. We want $x^2-y^2=a$ and $2xy=b$. Taking moduli of $(x+iy)^2=a+ib$ gives the third equation $x^2+y^2=|a+ib|$. Add and subtract:
\[
x^2=\frac{|a+ib|+a}2,\qquad y^2=\frac{|a+ib|-a}2.
\]
If $b>0$, then $x$ and $y$ have the same sign. If $b<0$, they have opposite signs. If $b=0$ and $a\ge0$, the roots are the ordinary real square roots $\pm\sqrt a$. If $b=0$ and $a<0$, the roots are $\pm i\sqrt{|a|}$. For instance $|3+4i|=5$, so $x^2=4$, $y^2=1$, and $b>0$: the square roots of $3+4i$ are $\pm(2+i)$. Check: $(2+i)^2=4+4i-1=3+4i$.

Every quadratic now has roots. Let $a\ne0$ and consider $az^2+bz+c=0$. Divide by $a$ and complete the square, exactly as over $\R$:
\[
\left(z+\frac{b}{2a}\right)^2=\frac{b^2-4ac}{4a^2}.
\]
The right side has a square root $\delta/(2a)$, where $\delta^2=b^2-4ac$, because every complex number has a square root. Therefore
\[
z=\frac{-b\pm\delta}{2a}.
\]
There are two roots, counted with multiplicity, unless you have chosen one of the two square roots and thrown the other away. Chapter~\ref{ch:polynomials} proves much more: every nonconstant polynomial with complex coefficients has a root.
\end{examplex}

\begin{examplex}[Bombelli's computation, checked]
Start from $2+i$, not from the cube root. Then
\[
(2+i)^2=4+4i+i^2=3+4i,
\]
and
\[
(2+i)^3=(3+4i)(2+i)=6+3i+8i+4i^2=2+11i.
\]
Since $\sqrt{-121}=11i$, this says $(2+i)^3=2+\sqrt{-121}$. Figure~\ref{fig:bombelli} draws the four points $1$, $2+i$, $(2+i)^2$ and $(2+i)^3$. Each step multiplies the modulus by $|2+i|=\sqrt5$ and adds the same angle, $\arctan(1/2)$. The points lie on a logarithmic spiral. The exponential form, in the next section, is the reason a repeated multiplication does that.

Bombelli's cube root was this number: $\sqrt[3]{2+\sqrt{-121}}=2+\sqrt{-1}$. Adding the conjugate gives $4$, which is the real root he wanted. The number $2+11i$ has two further cube roots. Exercise~\ref{ex:cardano} pairs each of them with a matching cube root of $2-11i$ and recovers $-2\pm\sqrt3$.
\end{examplex}

\begin{figure}[ht]
\centering
\fig{ch02-bombelli}
\caption{The powers $1$, $2+i$, $3+4i$ and $2+11i$. Each multiplication by $2+i$ multiplies lengths by $\sqrt5$ and adds the same angle, about $26.6^\circ$. The points sit on a logarithmic spiral.}
\label{fig:bombelli}
\end{figure}

\begin{think}
Multiply $1+i$ by itself. You get $2i$. Draw $1+i$ and $2i$ as arrows from the origin. What happened to the length, and what happened to the angle with the real axis? Now multiply $1+i$ by $i$, and look at $-1+i$. Guess the rule for a general product before you read the next section.
\end{think}

\subsection{Series with complex terms}
A complex sequence is two real sequences. Write $z_n=x_n+iy_n$ and $z=x+iy$. Then $z_n\to z$ if and only if $x_n\to x$ and $y_n\to y$. One direction uses $|x_n-x|\le|z_n-z|$ and $|y_n-y|\le|z_n-z|$. The other uses the triangle inequality,
\[
|z_n-z|\le|x_n-x|+|y_n-y|.
\]
Cauchy's criterion, absolute convergence, the Cauchy product and Dirichlet's test, all from Chapter~\ref{ch:sequences}, therefore extend to complex terms. For Cauchy's criterion, apply the real criterion to each coordinate. The other proofs are the same proofs, with the complex triangle inequality in place of the real one. Nothing in this chapter needs a complex derivative.

\section{The exponential form}

Addition is easy in coordinates. Multiplication is easy in modulus and argument. You already saw the pattern: $(1+i)^2=2i$. The length $\sqrt2$ was squared, and the angle $45^\circ$ was doubled, which lands on the imaginary axis. The exponential form is the dictionary that makes this a theorem. Figure~\ref{fig:arguments} is the one ambiguity in the dictionary. An argument is defined only up to a full turn.

\begin{figure}[ht]
\centering
\fig{ch02-arguments}
\caption{One point, many arguments. The short arc is an argument $\theta$. The dashed arc is $\theta+2\pi$. Both end at the same point, so a chosen argument is a representative of a class modulo $2\pi$.}
\label{fig:arguments}
\end{figure}

Use the real functions $\exp$, $\cos$ and $\sin$, and the constant $\pi$, with the properties listed in the foundational contract and compared in Proposition~\ref{prop:sincmp}. For $z=x+iy$, define
\[
\exp z=e^x(\cos y+i\sin y).
\]
When $y=0$ this is the real exponential. We write $e^z$ for $\exp z$.

Here is the calculation that makes the definition worth having. Let $z=x+iy$ and $w=x'+iy'$. The real addition formulas for cosine and sine, followed by the ordinary product in $\C$, give
\begin{align*}
\exp(z+w)&=e^{x+x'}\bigl(\cos(y+y')+i\sin(y+y')\bigr)\\
&=e^xe^{x'}\bigl((\cos y\cos y'-\sin y\sin y')+i(\sin y\cos y'+\cos y\sin y')\bigr)\\
&=e^x(\cos y+i\sin y)\cdot e^{x'}(\cos y'+i\sin y')\\
&=\exp z\cdot\exp w.
\end{align*}
Conjugation is the same computation with the sign of the imaginary part flipped: cosine is even and sine is odd, so $\overline{\exp z}=\exp\bar z$.

The series $\sum_{n\ge0}z^n/n!$ agrees with this definition. The identification uses Taylor's formula with remainder, and it is Exercise~\ref{ex:taylorln}(c) in Chapter~\ref{ch:calculus}. It is not an assumption of the algebra in this chapter.

\begin{proposition}[exponential form]\label{prop:expform}
\begin{enumerate}[label=(\alph*)]
\item $|e^{i\theta}|=1$, $e^{i\alpha}e^{i\beta}=e^{i(\alpha+\beta)}$, and hence, for every $n\in\Z$,
\[
(\cos\theta+i\sin\theta)^n=\cos(n\theta)+i\sin(n\theta).
\]
The last identity is de~Moivre's formula.
\item Every $z\in\C^*$ can be written $z=re^{i\theta}$ with $r=|z|$ and $\theta\in\R$. The real number $r$ is unique. The angle $\theta$ is unique modulo $2\pi$.
\item $\displaystyle\cos\theta=\frac{e^{i\theta}+e^{-i\theta}}2$ and $\displaystyle\sin\theta=\frac{e^{i\theta}-e^{-i\theta}}{2i}$.
\end{enumerate}
\end{proposition}
\begin{proof}
(a) By the product rule just proved, and by $\overline{e^{i\theta}}=e^{-i\theta}$,
\[
|e^{i\theta}|^2=e^{i\theta}e^{-i\theta}=e^0=1.
\]
The product $e^{i\alpha}e^{i\beta}=e^{i(\alpha+\beta)}$ is the same rule. De~Moivre's formula for $n\ge0$ is induction on that product. The case $n=0$ is $1$. If it holds for some $n\ge0$, multiply by $e^{i\theta}$. For $n<0$, use $e^{-i\theta}=(e^{i\theta})^{-1}$, which is legitimate because $|e^{i\theta}|=1$.

(c) From $\overline{e^{i\theta}}=e^{-i\theta}$ we have $e^{-i\theta}=\cos\theta-i\sin\theta$. Add that to $e^{i\theta}$ and divide by $2$. Subtract and divide by $2i$.

(b) \emph{Existence.} Put $r=|z|$ and $u=z/r=x+iy$, so $x^2+y^2=1$ and $x\in[-1,1]$. Cosine is continuous, $\cos0=1$ and $\cos\pi=-1$. The intermediate value theorem, Theorem~\ref{thm:ivt}, gives $\theta_0\in[0,\pi]$ with $\cos\theta_0=x$. Then $\sin^2\theta_0=1-x^2=y^2$. Proposition~\ref{prop:sincmp} says $\sin\ge0$ on $[0,\pi]$, so $\sin\theta_0=|y|$. If $y\ge0$, then $u=e^{i\theta_0}$. If $y<0$, then $u=\cos\theta_0-i\sin\theta_0=e^{-i\theta_0}$.

\emph{When is the cosine equal to $1$?} We claim that $\cos$ is strictly decreasing on $[0,\pi]$, and that the only $t\in[0,2\pi)$ with $\cos t=1$ is $t=0$. The contract gives $\sin0=\sin\pi=0$, $\sin(\pi/2)=1$, $\cos0=1$, $\cos(\pi/2)=0$, $\cos\pi=-1$, the period $2\pi$, and the two leastness statements used in Proposition~\ref{prop:sincmp}: $\pi/2$ is the least positive zero of cosine, and $\pi$ is the least positive angle at which cosine equals $-1$.

On $[0,\pi/2)$, cosine is positive: it starts at $1$, and a negative value would force a zero before $\pi/2$, contradicting the leastness of that zero. Thus $\sin'=\cos>0$. The mean value theorem, Theorem~\ref{thm:mvt}, says that sine is strictly increasing on $[0,\pi/2]$, from $0$ to $1$. In particular $\sin x>0$ on $(0,\pi/2]$. On $[0,\pi]$ we already know $\sin\ge0$, so $\cos'=-\sin\le0$ and cosine is non-increasing there. Suppose $\sin x_0=0$ for some $x_0\in(\pi/2,\pi)$. Then $\cos x_0=\pm1$. The value $+1$ is impossible, because cosine is non-increasing and $\cos(\pi/2)=0$. The value $-1$ is impossible, because $\pi$ is the least positive angle where cosine equals $-1$. Thus $\sin>0$ on $(\pi/2,\pi)$ as well.

Now take $0\le a<b\le\pi$. The mean value theorem gives $\cos b-\cos a=-\sin c\cdot(b-a)$ for some $c\in(a,b)$. The sine on the right is positive, so $\cos b<\cos a$. Cosine is strictly decreasing on $[0,\pi]$, hence injective there, and $\cos t=1$ on that interval only at $t=0$. If instead $t\in(\pi,2\pi)$, then $2\pi-t\in(0,\pi)$ and $\cos t=\cos(2\pi-t)<1$.

\emph{Uniqueness.} Suppose $e^{i\theta}=e^{i\theta'}$. Then $e^{i(\theta-\theta')}=1$, so $\cos(\theta-\theta')=1$ and $\sin(\theta-\theta')=0$. Write $\theta-\theta'=2\pi k+t$ with $k\in\Z$ and $t\in[0,2\pi)$. Periodicity gives $e^{it}=1$, so $\cos t=1$. The previous paragraph forces $t=0$. Thus $\theta-\theta'\in2\pi\Z$.
\end{proof}

\portrait{demoivre}{Abraham de Moivre (1667--1754). Portrait, 1736, artist unknown. Public domain, via Wikimedia Commons.}

The number $\theta$ in $z=re^{i\theta}$ is an \emph{argument} of $z$. Statement (b) is the geometry promised by the pause above. If $z=re^{i\theta}$ and $w=se^{i\varphi}$, then
\[
zw=rs\,e^{i(\theta+\varphi)}.
\]
Multiplying by $w$ stretches every arrow by $s$ and turns it by $\varphi$. Figure~\ref{fig:multiply} is one case. Multiplying by $i=e^{i\pi/2}$ is a quarter turn anticlockwise. Two quarter turns are a half turn, and a half turn sends $1$ to $-1$. That is $i^2=-1$, drawn as a walk around the unit circle. Figure~\ref{fig:turn} is the walk. There are two quarter turns that square to a half turn, anticlockwise and clockwise, and those are the two square roots of $-1$, namely $i$ and $-i$.

The same dictionary checks the square root of $i$ from the previous section. Since $i=e^{i\pi/2}$, one square root is $e^{i\pi/4}=\cos(\pi/4)+i\sin(\pi/4)=(1+i)/\sqrt2$. The other is $e^{i(\pi/4+\pi)}=e^{i5\pi/4}=-(1+i)/\sqrt2$. The angle was divided by $2$. Adding $\pi$ gives the second root, because $e^{i\pi}=-1$ and $(-1)^2=1$.

Bombelli's spiral is the same rule. The number $2+i$ has modulus $\sqrt5$ and an argument $\theta=\arctan(1/2)$. Its powers are $(\sqrt5)^n e^{in\theta}$.

\begin{figure}[ht]
\centering
\fig{ch02-multiply}
\caption{Multiplication by $w=1.2\,e^{i\cdot50^\circ}$ stretches $z=2+i$ by $1.2$ and turns it by $50^\circ$. In general $|zw|=|z|\,|w|$, and the arguments add.}
\label{fig:multiply}
\end{figure}

\begin{figure}[ht]
\centering
\fig{ch02-turn}
\caption{Multiplication by $i$ is a quarter turn anticlockwise. The first turn sends $1$ to $i$. The second sends $i$ to $-1$. Two quarter turns are a half turn, which is $i^2=-1$.}
\label{fig:turn}
\end{figure}

Two habits follow, and both will be used again.

To \emph{linearise} a power of cosine or sine, write it through part (c) and expand with the binomial theorem. The powers $e^{ik\theta}$ come back to cosines and sines of multiple angles.

To \emph{add} two turns of the same length, factor the half-angle:
\begin{align*}
e^{ia}+e^{ib}&=e^{i(a+b)/2}\left(e^{i(a-b)/2}+e^{-i(a-b)/2}\right)\eqbr=2\cos\frac{a-b}2\;e^{i(a+b)/2}.
\end{align*}
In particular,
\[
1+e^{i\theta}=2\cos\frac\theta2\,e^{i\theta/2},\qquad 1-e^{i\theta}=-2i\sin\frac\theta2\,e^{i\theta/2},
\]
and therefore $|1-e^{i\theta}|=2\bigl|\sin(\theta/2)\bigr|$. Figure~\ref{fig:rhombus} is the same identity without formulas. Two arrows of equal length form a rhombus. Their sum runs along the bisector, at angle $(a+b)/2$, and its length is $2\bigl|\cos((a-b)/2)\bigr|$. If that cosine is negative, the sum points opposite the bisector you chose. The algebra is that picture.

\begin{figure}[ht]
\centering
\fig{ch02-rhombus}
\caption{The sum of two unit arrows points along the bisector of the angle between them. The dashed sides complete the rhombus. The length of the sum is $2\bigl|\cos((a-b)/2)\bigr|$.}
\label{fig:rhombus}
\end{figure}

\begin{example}
Linearise a cube. With part (c) of Proposition~\ref{prop:expform},
\begin{align*}
\cos^3\theta&=\frac18\bigl(e^{i\theta}+e^{-i\theta}\bigr)^3\eqbr=\frac18\bigl(e^{3i\theta}+3e^{i\theta}+3e^{-i\theta}+e^{-3i\theta}\bigr)\eqbr=\frac{\cos3\theta+3\cos\theta}4.
\end{align*}
At $\theta=0$ both sides are $1$. That is a cheap check, and it catches a lost factor of $2$.
\end{example}

\begin{examplex}[the addition formulas, as a check]
Take real and imaginary parts of $e^{i(\alpha+\beta)}=e^{i\alpha}e^{i\beta}$:
\begin{align*}
\cos(\alpha+\beta)&=\cos\alpha\cos\beta-\sin\alpha\sin\beta,\\
\sin(\alpha+\beta)&=\sin\alpha\cos\beta+\cos\alpha\sin\beta.
\end{align*}
These are the real addition formulas that went into the definition of $\exp$. Recovering them is a consistency check, not an independent proof. The double-angle formulas are the case $\alpha=\beta$. The identity $\cos2\theta=2\cos^2\theta-1$ is the real part of $(e^{i\theta})^2$, together with $\sin^2\theta=1-\cos^2\theta$. The sum-to-product formulas are the half-angle factorisation above, read in the other direction.
\end{examplex}

\section{Workshop: a turn is not a place}

A complex number does two jobs. Added, it is a position. Multiplied, it is an instruction: stretch by this factor, turn by this angle. The next example is what happens if you obey the instruction and forget that the lengths are changing.

\begin{examplex}[a robot that does not come home]
A robot starts at $0$ and takes steps $1,q,q^2,q^3$, in that order. Its position after all four steps is the sum $1+q+q^2+q^3$.

If $q=i$, the steps are $1$, $i$, $-1$, $-i$. The positions are $0$, $1$, $1+i$, $i$, $0$. It walked a square and came home. If $q=(1+i)/\sqrt2$, then $q^8=1$ and eight unit steps close. The reason is the geometric sum below.

If $q=1+i$, each step turns by $45^\circ$ and is $\sqrt2$ times as long as the one before. Four equal steps of $45^\circ$ do not close. The square is the path for $q=i$, a turn of $90^\circ$. After four steps the position is
\[
1+(1+i)+(1+i)^2+(1+i)^3=1+(1+i)+2i+(-2+2i)=5i,
\]
not $0$. Figure~\ref{fig:robot} draws the path. The turn stays $45^\circ$, and the lengths grow, so the path swells and stops at $5i$.

In general, if $q\ne1$, the position after $n$ steps is $(1-q^n)/(1-q)$. If $q=1$, every step is $1$ and the position is $n$. Separate that case before you divide.
\end{examplex}

\begin{figure}[ht]
\centering
\fig{review02-robot}
\caption{Steps with multiplier $1+i$. The positions are $0$, $1$, $2+i$, $2+3i$ and $5i$. Each arrow is a step, not the position vector. The turn is $45^\circ$ every time, and the lengths grow by $\sqrt2$.}
\label{fig:robot}
\end{figure}

\begin{examplex}[three cube roots, and an argument that jumps]
Solve $z^3=8i$. First write the right side in exponential form, and write every argument, not only the one that sits in $(-\pi,\pi]$. Since $8i$ has modulus $8$ and one argument $\pi/2$,
\[
8i=8e^{i(\pi/2+2\pi k)},\qquad k\in\Z.
\]
A cube root has modulus $8^{1/3}=2$, and an argument $(\pi/2+2\pi k)/3=\pi/6+2\pi k/3$. The three values $k=0,1,2$ give three angles in an interval of length $2\pi$, namely $\pi/6$, $5\pi/6$ and $3\pi/2$. The roots are
\[
2e^{i\pi/6}=\sqrt3+i,\qquad 2e^{i5\pi/6}=-\sqrt3+i,\qquad 2e^{i3\pi/2}=-2i.
\]
Cubing any one of them returns $8i$. Their arguments differ by less than a full turn, so they are distinct. Part (b) of Proposition~\ref{prop:expform} says there are no others: every nonzero number has a modulus and an argument class, and dividing the class by $3$ produces exactly three classes.

The same arithmetic shows why a chosen argument of a product can jump by $2\pi$. Take principal arguments in $(-\pi,\pi]$. The number $e^{i3\pi/4}$ has principal argument $3\pi/4$. The product of two copies has argument $3\pi/2$ if you simply add, and $3\pi/2$ is not in $(-\pi,\pi]$. The representative in that interval is $3\pi/2-2\pi=-\pi/2$. Arguments add as classes modulo $2\pi$. Chosen representatives do not always add as real numbers.
\end{examplex}

\begin{reminder}[a square-root symbol has to choose]
The equation $w^2=z$, for $z\ne0$, has two solutions. A square-root \emph{function} returns one of them. The real identity $\sqrt{ab}=\sqrt a\sqrt b$ can fail for that choice. For the principal square root, $\sqrt{-1}=i$, so $\sqrt{-1}\,\sqrt{-1}=-1$, while $\sqrt{(-1)(-1)}=\sqrt1=1$. The symbol had to pick one root, and the two picks did not multiply the way the two sets of roots would. An equation that lists every root is safer than an unannounced branch.
\end{reminder}

\section{Roots of unity}

\begin{think}
Draw the solutions of $z^6=1$ in the plane. Which of them have all six solutions among their powers? Which visit only a smaller regular polygon? What is the sum of all six? Answer from the picture before you compute.
\end{think}

\begin{proposition}[roots of unity]\label{prop:roots}
Let $n\ge1$ and $\omega=e^{2i\pi/n}$.
\begin{enumerate}[label=(\alph*)]
\item The equation $z^n=1$ has exactly $n$ solutions, namely $\omega^k$ for $0\le k\le n-1$. They form a group $\mathbb U_n$ under multiplication, and
\[
X^n-1=\prod_{k=0}^{n-1}\bigl(X-\omega^k\bigr).
\]
\item The order of $\omega^k$, meaning the least positive integer $m$ with $(\omega^k)^m=1$, is $n/\gcd(k,n)$. In particular $\omega^k$ is \emph{primitive} --- its powers are all of $\mathbb U_n$ --- exactly when $\gcd(k,n)=1$.
\end{enumerate}
\end{proposition}
\begin{proof}
(a) If $z^n=1$, then $|z|^n=1$, so $|z|=1$. Write $z=e^{i\theta}$. Then $e^{in\theta}=1$. By the uniqueness in Proposition~\ref{prop:expform}, $n\theta\in2\pi\Z$, so $\theta=2\pi k/n$ for some $k\in\Z$, and $k$ matters only modulo $n$. Thus $z=\omega^k$ for a unique $k$ with $0\le k\le n-1$. Each such power is a solution, because $(\omega^k)^n=\omega^{kn}=(e^{2i\pi})^k=1$. The $n$ arguments $2\pi k/n$ lie in $[0,2\pi)$ and are distinct, so the $n$ points are distinct.

The product of two solutions is a solution, because $(zw)^n=z^nw^n=1$. The inverse of a solution is a solution. The identity is $1=\omega^0$. So $\mathbb U_n$ is a group.

For the polynomial, fix $k$ and use $(\omega^k)^n=1$:
\[
X^n-1=X^n-(\omega^k)^n=(X-\omega^k)\sum_{l=0}^{n-1}X^{n-1-l}(\omega^k)^l.
\]
The quotient is monic of degree $n-1$. Every $\omega^j$ with $j\not\equiv k\pmod n$ is a root of $X^n-1$ and is not equal to $\omega^k$, so it is a root of the quotient. Remove one linear factor at a time. After $n$ steps, $X^n-1$ equals the product of the $n$ distinct factors $X-\omega^k$, up to a constant, and both sides are monic, so the constant is $1$. Proposition~\ref{prop:rootsfactor} is this argument for a general polynomial.

(b) Let $d=\gcd(k,n)$, and write $k=dk'$ and $n=dn'$, so $\gcd(k',n')=1$. Then $(\omega^k)^m=\omega^{km}=1$ if and only if $n$ divides $km$, because $\omega^t=1$ if and only if $n$ divides $t$. That divisibility is $dn'\mid dk'm$, hence $n'\mid k'm$. Euclid's lemma, from the sums chapter, says that a divisor of $k'm$ which is coprime to $k'$ must divide $m$. Thus $n'\mid m$. The least positive such $m$ is $n'=n/d$.

The powers of $\omega^k$ are $n$ distinct elements exactly when the order is $n$, which is exactly when $d=1$.
\end{proof}

The $n$th roots of unity are the vertices of the regular $n$-gon inscribed in the unit circle, with one vertex at $1$. The order of a root is the number of distinct powers it has. A primitive root visits every vertex before it returns to $1$. A root that is not primitive visits a smaller regular polygon, the subgroup of roots of unity of its own order.

For $n=6$, let $\omega=e^{2i\pi/6}=e^{i\pi/3}$. Then $\omega$ and $\omega^5$ are primitive. The powers $\omega^2$ and $\omega^4$ have order $3$: they are the non-real cube roots of unity, and together with $1$ they form $\mathbb U_3$. The power $\omega^3=-1$ has order $2$. Figure~\ref{fig:orders} is that list. Figure~\ref{fig:roots} is the same idea for $n=12$, together with the chords that will give the product of sines.

\begin{figure}[ht]
\centering
\fig{ch02-orders}
\caption{The sixth roots of unity. Rust marks the primitive roots $\omega$ and $\omega^5$, the ones whose powers visit every vertex. Ochre marks $\omega^2$ and $\omega^4$, which have order $3$. The dashed triangle is the subgroup $\mathbb U_3$. The point $-1$ has order $2$.}
\label{fig:orders}
\end{figure}

\begin{figure}[ht]
\centering
\fig{ch02-roots}
\caption{Left: the twelve twelfth roots of unity. The primitive ones, $\omega^k$ with $\gcd(k,12)=1$, are in rust. The others generate $\mathbb U_d$ for a proper divisor $d$ of $12$. Right: the six chords from $1$ in the regular heptagon. Each has length $|1-\omega^k|=2\sin(k\pi/7)$, and the product of the six lengths is $7$.}
\label{fig:roots}
\end{figure}

\begin{proposition}[orthogonality]\label{prop:orth}
Let $\omega=e^{2i\pi/n}$ and $j\in\Z$. Then
\[
\sum_{k=0}^{n-1}\omega^{jk}=\begin{cases}n&\text{if }n\mid j,\\0&\text{otherwise.}\end{cases}
\]
\end{proposition}
\begin{proof}
If $n$ divides $j$, then every term is $\omega^{j}=1$, and the sum is $n$. If not, then $q=\omega^j\ne1$, while $q^n=(\omega^n)^j=1$. The geometric sum is $(q^n-1)/(q-1)=0$.
\end{proof}

The case $j=1$ is the sum of the vertices. It is $0$. The centre of mass of a regular polygon inscribed in the unit circle, with one vertex at $1$, is the origin. That is the answer to the pause: the six sixth roots add to $0$, and you can see it, because they come in three opposite pairs. The name orthogonality comes from a second reading of the same sum. The vectors
\[
v_j=\bigl(1,\omega^j,\omega^{2j},\dots,\omega^{(n-1)j}\bigr),\qquad j=0,\dots,n-1,
\]
are pairwise orthogonal for the Hermitian product $\langle v,w\rangle=\sum_k v_k\overline{w_k}$. Indeed $\langle v_j,v_l\rangle=\sum_k\omega^{(j-l)k}$, which is $0$ when $j\not\equiv l\pmod n$. These vectors are the Fourier basis of $\C^n$. Exercise~\ref{ex:circulant} puts them to work as eigenvectors.

\begin{examplex}[the cosine of $2\pi/5$]
Let $\omega=e^{2i\pi/5}$, and set $c=\omega+\omega^4$ and $d=\omega^2+\omega^3$. The sum of all five fifth roots is $0$, by Proposition~\ref{prop:orth}, so
\[
c+d=-1.
\]
For the product, reduce exponents modulo $5$, using $\omega^5=1$:
\begin{align*}
cd&=(\omega+\omega^4)(\omega^2+\omega^3)=\omega^3+\omega^4+\omega^6+\omega^7\eqbr=\omega^3+\omega^4+\omega+\omega^2=-1.
\end{align*}
Thus $c$ and $d$ are the two roots of $T^2+T-1=0$, namely $(-1\pm\sqrt5)/2$.

Which is which? Since $\omega^4=\bar\omega$, the sum $c$ is real and equals $2\cos(2\pi/5)$. The angle $2\pi/5$ lies in $(0,\pi/2)$, so the cosine is positive, and
\[
c=\frac{-1+\sqrt5}2,\qquad\cos\frac{2\pi}5=\frac{\sqrt5-1}4.
\]
The other root is negative: $d=(-1-\sqrt5)/2$. The same pairing gives $d=2\cos(4\pi/5)$. Since $\cos(\pi/5)=-\cos(4\pi/5)$,
\[
\cos\frac\pi5=-\frac d2=\frac{\sqrt5+1}4.
\]
Figure~\ref{fig:pentagon} draws the two conjugate pairs and the two real sums. Both values are built from rational operations and a single square root. That is why a regular pentagon can be drawn with ruler and compass. The regular heptagon cannot, and the reason is in the challenges, with the proof in Volume~3.
\end{examplex}

\begin{figure}[ht]
\centering
\fig{ch02-pentagon}
\caption{The fifth roots of unity. The rust chord joins $\omega$ to $\omega^4$, and their vector sum is the positive real number $c=2\cos(2\pi/5)$. The ochre chord joins $\omega^2$ to $\omega^3$, and their sum is the negative real number $d=2\cos(4\pi/5)$.}
\label{fig:pentagon}
\end{figure}

\begin{think}
Inscribe a regular heptagon in the unit circle, and multiply the lengths of the six chords from one vertex to the other six (Figure~\ref{fig:roots}, right). The product is an integer. Which integer would you guess for a regular $n$-gon?
\end{think}

\begin{examplex}[the product of sines]
The guess from the picture is $n$. Here is the proof. From Proposition~\ref{prop:roots},
\[
\prod_{k=1}^{n-1}\bigl(X-\omega^k\bigr)=\frac{X^n-1}{X-1}=1+X+\dots+X^{n-1},
\]
for $n\ge2$. Set $X=1$. The right side is $n$, so $\prod_{k=1}^{n-1}(1-\omega^k)=n$. Take moduli, and use $|1-e^{i\theta}|=2\bigl|\sin(\theta/2)\bigr|$ with $\theta=2\pi k/n$:
\[
n=\prod_{k=1}^{n-1}2\sin\frac{k\pi}n=2^{n-1}\prod_{k=1}^{n-1}\sin\frac{k\pi}n.
\]
Every sine is positive, because $0<k\pi/n<\pi$. Therefore
\[
\prod_{k=1}^{n-1}\sin\frac{k\pi}n=\frac n{2^{n-1}}.
\]
For $n=7$ the right side is $7/64$, which is the product in the epigraph. The six chord lengths multiply to $7$, not to $7/64$: each chord contributes a factor $2$, and $2^6=64$.
\end{examplex}

\section{The roots-of-unity filter}

\begin{think}
The sum of all binomial coefficients $\sum_k\binom nk$ is $2^n$, which is $(1+1)^n$. The alternating sum $\sum_k(-1)^k\binom nk$ is $0$ for $n\ge1$, which is $(1-1)^n$. Adding those two evaluations and dividing by $2$ keeps the terms with even $k$ and kills the terms with odd $k$. You want only the indices divisible by $3$. What should replace $1$ and $-1$?
\end{think}

Proposition~\ref{prop:orth} extracts, from a polynomial, the terms whose exponent lies in one residue class.

\begin{theorem}[roots-of-unity filter]\label{thm:filter}
Let $P(X)=\sum_k a_kX^k$, let $m\ge1$, let $\omega=e^{2i\pi/m}$ and let $r\in\Z$. Then
\[
\sum_{k\equiv r\pmod m}a_k=\frac1m\sum_{j=0}^{m-1}\omega^{-jr}P(\omega^j).
\]
\end{theorem}
\begin{proof}
Substitute the sum for $P$. The right side is
\[
\sum_k a_k\cdot\frac1m\sum_{j=0}^{m-1}\omega^{j(k-r)}.
\]
By Proposition~\ref{prop:orth}, the inner average is $1$ when $m$ divides $k-r$, and $0$ otherwise.
\end{proof}

\begin{figure}[ht]
\centering
\fig{ch02-filter}
\caption{The filter for $m=3$. If $3$ does not divide $k$, the three numbers $1$, $\omega^k$ and $\omega^{2k}$ are the vertices of an equilateral triangle, and they add to $0$. If $3$ divides $k$, all three equal $1$, and they add to $3$. Averaging $P(\omega^j)$ therefore keeps the coefficients $a_k$ with $3\mid k$ and drops the rest.}
\label{fig:filter}
\end{figure}

\begin{example}
Take $P=(1+X)^n$, $m=2$ and $r=0$. Then $\omega=-1$, and
\[
\sum_k\binom n{2k}=\frac12\bigl((1+1)^n+(1-1)^n\bigr)=\frac12\bigl(2^n+0^n\bigr).
\]
For $n\ge1$ the second term vanishes, and the sum is $2^{n-1}$. For $n=0$ the sum is $1$, and $0^0$ in this formula is the empty product, which we read as $1$. The case $m=3$ is Exercise~\ref{ex:mod3}. It counts the subsets whose size is a multiple of $3$.
\end{example}

\section{Geometry with complex numbers}

Every map $z\mapsto az+b$ with $a\ne0$ is a \emph{direct similarity}. Write $a=re^{i\theta}$. The map turns by $\theta$, scales by $r$, and translates by $b$, in that order. It preserves angles and ratios of lengths, so it sends equilateral triangles to equilateral triangles and squares to squares. A geometric statement that is unchanged by similarities may be proved after you move two convenient points to $0$ and $1$.

The indirect similarities are $z\mapsto a\bar z+b$ with $a\ne0$. We do not need them. The criterion below handles most problems about equilateral triangles.

Write $j=e^{2i\pi/3}=(-1+i\sqrt3)/2$. Then $j^3=1$, $1+j+j^2=0$, and $\bar j=j^2$. Also $-j^2=e^{i\pi/3}$, which is the rotation by $60^\circ$ anticlockwise. Check the last one: $-j^2=e^{i\pi}e^{-2i\pi/3}=e^{i\pi/3}$.

\begin{proposition}[equilateral triangles]\label{prop:equilateral}
For $a,b,c\in\C$, the following are equivalent:
\begin{enumerate}[label=(\roman*)]
\item the triangle $abc$ is equilateral, allowing the degenerate triangle $a=b=c$;
\item $a+jb+j^2c=0$ or $a+j^2b+jc=0$;
\item $a^2+b^2+c^2=ab+bc+ca$.
\end{enumerate}
The first alternative in (ii) holds when $a,b,c$ are in anticlockwise order. The second holds when they are in clockwise order.
\end{proposition}
\begin{proof}
The triangle is equilateral and anticlockwise exactly when $c-a$ is the rotate of $b-a$ by $60^\circ$:
\[
c-a=-j^2(b-a).
\]
Expand: $c-a=-j^2b+j^2a$, so $c+j^2b-a(1+j^2)=0$. Since $1+j+j^2=0$, we have $-(1+j^2)=j$, and the equation is $ja+j^2b+c=0$. Multiply by $j^2$. Using $j^3=1$ and $j^4=j$, you get $a+jb+j^2c=0$.

The clockwise case is the rotation by $-60^\circ$, namely $c-a=-j(b-a)$. The same algebra produces $a+j^2b+jc=0$.

For (iii), expand the product, using $j^3=1$ and $j+j^2=-1$. The cross terms are $ab(j+j^2)$, $bc(j+j^2)$ and $ca(j+j^2)$, so
\[
(a+jb+j^2c)(a+j^2b+jc)=a^2+b^2+c^2-ab-bc-ca.
\]
The product is zero exactly when one of the factors is zero, because $\C$ is a field.
\end{proof}

\begin{figure}[ht]
\centering
\fig{ch02-napoleon}
\caption{On the sides of any triangle, build equilateral triangles outward. Their centres, the ochre dots, are themselves the vertices of an equilateral triangle. Exercise~\ref{ex:napoleon} is the proof. The dashed triangle does not depend on the shape you started with.}
\label{fig:napoleon}
\end{figure}

\section{Cyclotomic polynomials}

The polynomial $X^n-1$ factors over $\C$ into the $n$ linear factors $X-\omega^k$. Group the roots by their order, and the factors that come out have integer coefficients. Those factors are the cyclotomic polynomials. Volume~2 uses their values at integers to produce primes congruent to $1$ modulo $n$. Volume~3 proves that each of them is irreducible over $\Q$.

\begin{definition}\label{def:cyclotomic}
The $n$th \emph{cyclotomic polynomial} is
\[
\Phi_n(X)=\prod(X-\zeta),
\]
the product over the primitive $n$th roots of unity, that is, over $\zeta=\omega^k$ with $0\le k\le n-1$ and $\gcd(k,n)=1$. Its degree is the number of such $k$. Exercise~\ref{ex:phi} in Chapter~\ref{ch:counting} computes that number: it is Euler's $\varphi(n)$.
\end{definition}

\begin{proposition}\label{prop:cyclotomic}
\begin{enumerate}[label=(\alph*)]
\item $X^n-1=\prod_{d\mid n}\Phi_d(X)$, the product over the positive divisors of $n$.
\item $\Phi_n$ is monic and has integer coefficients.
\item If $p$ is prime, then $\Phi_p(X)=1+X+\dots+X^{p-1}$ and $\Phi_{p^k}(X)=\Phi_p\bigl(X^{p^{k-1}}\bigr)$.
\end{enumerate}
\end{proposition}
\begin{proof}
(a) Every $n$th root of unity has some order $d$, and Proposition~\ref{prop:roots}(b) says that $d$ divides $n$. The roots of order exactly $d$ are the primitive $d$th roots of unity. They are $n$th roots of unity because $d$ divides $n$. The roots of $X^n-1$ are simple, and they are partitioned by their orders, so
\[
X^n-1=\prod_{d\mid n}\Phi_d(X).
\]

(b) Proceed by strong induction on $n$. The polynomial $\Phi_1=X-1$ is monic with integer coefficients. For $n>1$, part (a) gives $X^n-1=\Phi_n\cdot Q$ with $Q=\prod_{d\mid n,\,d<n}\Phi_d$. By the induction hypothesis, $Q$ is monic and lies in $\Z[X]$.

Divide $X^n-1$ by $Q$ by the usual algorithm, subtracting at each step an integer multiple of a power of $X$ times $Q$. The leading coefficient of $Q$ is $1$, so that multiplier is the current leading coefficient, which starts in $\Z$ and stays in $\Z$. Stop when the remainder $R$ has degree less than $\deg Q$. This produces $X^n-1=QS+R$ with $S,R\in\Z[X]$.

The same quotient and remainder are unique even over $\C$. If $X^n-1=QS_1+R_1=QS_2+R_2$ with both remainders of degree less than $\deg Q$, then $Q(S_1-S_2)=R_2-R_1$. If $S_1\ne S_2$, the left side has degree at least $\deg Q$ and the right side has degree less than $\deg Q$. So $S_1=S_2$ and $R_1=R_2$. Over $\C$ we already know that the quotient is $\Phi_n$ and the remainder is $0$. Therefore $\Phi_n=S$, which is monic and lies in $\Z[X]$. Proposition~\ref{prop:polydiv} is the general division algorithm.

(c) Every $p$th root of unity except $1$ is primitive, because the order divides $p$ and is not $1$. Thus $\Phi_p(X)=(X^p-1)/(X-1)$. For $p^k$, the $p^k$th roots that are not primitive are the roots whose order divides $p^{k-1}$, and those are the roots of $X^{p^{k-1}}-1$. Therefore
\[
\Phi_{p^k}(X)=\frac{X^{p^k}-1}{X^{p^{k-1}}-1}=\Phi_p\bigl(X^{p^{k-1}}\bigr).
\]
\end{proof}

\begin{examplex}[a division, and a false guess]
The factorisation for $n=6$ can be written by hand. The positive divisors are $1,2,3,6$, and
\[
\Phi_1=X-1,\qquad\Phi_2=X+1,\qquad\Phi_3=X^2+X+1.
\]
Their product is $(X^2-1)(X^2+X+1)=X^4+X^3-X-1$. Divide $X^6-1$ by that monic polynomial. The first subtraction uses $X^2$:
\[
(X^6-1)-X^2(X^4+X^3-X-1)=-X^5+X^3+X^2-1.
\]
The next subtraction uses $-X$:
\[
(-X^5+X^3+X^2-1)-(-X)(X^4+X^3-X-1)=X^4+X^3-X-1.
\]
One further subtraction, of the divisor itself, leaves remainder $0$. The quotient is $X^2-X+1$. So $\Phi_6=X^2-X+1$.

A few more, from the same recipe: $\Phi_4=X^2+1$, $\Phi_5=X^4+X^3+X^2+X+1$, $\Phi_8=X^4+1$, $\Phi_{10}=X^4-X^3+X^2-X+1$, $\Phi_{12}=X^4-X^2+1$. For every $n\le104$, every coefficient is $0$, $1$, or $-1$. It is tempting to guess that this always happens. It does not. The polynomial $\Phi_{105}$ has exactly two coefficients equal to $-2$, those of $X^7$ and of $X^{41}$. Migotti recorded a coefficient $-2$ in $1883$. In a letter to Landau in $1931$, Schur showed that the coefficients are unbounded as $n$ varies. Emma Lehmer published his argument in $1936$. Exercise~\ref{ex:cyclo} computes a few more polynomials. Chapter~\ref{ch:polynomials} proves that $\Phi_p$ is irreducible over $\Q$. The same fact for every $n$ is proved in Volume~3. Gauss treated the prime case in the \emph{Disquisitiones}; the argument for general $n$ came later.
\end{examplex}

\portrait{gauss}{Carl Friedrich Gauss (1777--1855). Portrait by Christian Albrecht Jensen, 1840. Public domain, via Wikimedia Commons.}

The $17$-gon is the same circle, one step further. A regular $p$-gon with $p$ prime is constructible with ruler and compass only if $p-1$ is a power of $2$, because the degree of $\Phi_p$ is $p-1$, and each ruler-and-compass step solves a quadratic. For $p=5$ that degree is $4=2^2$, and the example above reached $\cos(2\pi/5)$ with one square root. For $p=17$ the degree is $16=2^4$, which is why Gauss's diary entry is possible. For $p=7$ the degree is $6$, which is not a power of $2$. Volume~3 proves the precise criterion. The challenges ask you to test it for $n\le30$, and to see the cubic that stops the heptagon.

\section*{What this chapter prepares}
A nonzero complex number is a length and a turn. The length multiplies, by Brahmagupta's identity. The turn adds, up to a full circle, which is why $z^n=1$ has $n$ solutions rather than one. Those solutions are the vertices of a regular $n$-gon. Their sum is $0$. The filter is that cancellation used as a sieve: average a polynomial against the $m$th roots of unity, and one residue class of coefficients remains. Grouping the vertices by order factors $X^n-1$ into cyclotomic polynomials with integer coefficients.

The next chapter asks how far factorisation goes for a general polynomial. Division produces a quotient and a remainder. A root produces a linear factor. The existence of a root, for every nonconstant polynomial with complex coefficients, uses the exponential form from this chapter: an $n$th root is a length, and an angle divided by $n$. One case is already proved here. Every complex number has a square root, so every quadratic equation has roots in $\C$.

Volume~2, \emph{The Integers: Primes, Patterns and Proofs}, uses a square of the Gauss sum from Exercise~\ref{ex:gauss} in a proof of quadratic reciprocity, and uses the values of $\Phi_n$ at integers to reach infinitely many primes congruent to $1$ modulo $n$. Volume~3 proves that every cyclotomic polynomial is irreducible, and decides which regular polygons can be drawn with ruler and compass. Volume~5 turns the filter into a fast way to multiply polynomials, and determines the sign of the Gauss sum.

\section{Exercises}

\begin{exercise}[$\star$]\label{ex:lin}
Linearise $\sin^4\theta$. After the integration tools of Chapter~\ref{ch:calculus}, deduce $\displaystyle\int_0^{\pi}\sin^4\theta\dd\theta$.
\end{exercise}

\begin{exercise}[$\star$]\label{ex:binomcos}
Compute $\displaystyle\sum_{k=0}^n\binom nk\cos(k\theta)$ in closed form.
\end{exercise}

\begin{exercise}[$\star\star$]\label{ex:dirichlet}
\emph{The Dirichlet kernel.} For $\theta\notin2\pi\Z$, show that
\[
D_n(\theta)=\sum_{k=-n}^ne^{ik\theta}=\frac{\sin\bigl((n+\tfrac12)\theta\bigr)}{\sin(\theta/2)},
\]
and deduce $\displaystyle\sum_{k=1}^n\cos(k\theta)$.
\end{exercise}

\begin{exercise}[$\star\star$]\label{ex:mod3}
\emph{Bridge: binomial coefficients and roots of unity.}
How many subsets of $\{1,\dots,n\}$ have size divisible by $3$?
\end{exercise}

\begin{exercise}[$\star\star$]\label{ex:sines}
Let $n\ge2$ and $\omega=e^{2i\pi/n}$. Show that $\prod_{k=1}^{n-1}(1+\omega^k)$ equals $1$ when $n$ is odd, and equals $0$ when $n$ is even.
\end{exercise}

\begin{exercise}[$\star\star$]\label{ex:pentagon}
The example in the text gives $\cos(2\pi/5)=(\sqrt5-1)/4$ and $\cos(\pi/5)=(\sqrt5+1)/4$.
\begin{enumerate}[label=(\alph*)]
\item Check the double-angle formula against these values:
\[
\cos\frac{4\pi}5=2\cos^2\frac{2\pi}5-1.
\]
\item Deduce $\sin(\pi/5)$ and $\sin(2\pi/5)$.
\end{enumerate}
The side of a regular pentagon inscribed in the unit circle is $2\sin(\pi/5)$.
\end{exercise}

\begin{exercise}[$\star\star\star$]\label{ex:circulant}
\emph{Bridge: roots of unity and linear algebra.}
\emph{Circulant determinants, to be read again after Chapter~\ref{ch:linear}.} Let $a_0,\dots,a_{n-1}\in\C$, let $P(X)=\sum_ka_kX^k$, and let $C$ be the $n\times n$ matrix with entries $C_{r,s}=a_{(s-r)\bmod n}$. Each row is the previous row, shifted right by one.
\begin{enumerate}[label=(\alph*)]
\item With $\omega=e^{2i\pi/n}$ and $v_j=\bigl(1,\omega^j,\omega^{2j},\dots,\omega^{(n-1)j}\bigr)^{\mathsf T}$, show that $Cv_j=P(\omega^j)\,v_j$.
\item Deduce that $\det C=\prod_{j=0}^{n-1}P(\omega^j)$.
\item For $n\ge2$, compute the determinant of the circulant matrix with first row $(1,1,0,\dots,0)$.
\end{enumerate}
\end{exercise}

\begin{exercise}[$\star\star\star$]\label{ex:gauss}
\emph{A quadratic Gauss sum.} Let $p$ be an odd prime, $\omega=e^{2i\pi/p}$ and $G=\sum_{a=0}^{p-1}\omega^{a^2}$. Show that $|G|^2=p$.
\end{exercise}

\subsection*{Further exercises}

\begin{exercise}[$\star$]\label{ex:sqrtz}
Find the square roots of $-7+24i$ and of $2i$, and solve $z^2-(1+i)z+i=0$.
\end{exercise}

\begin{exercise}[$\star$]\label{ex:sumdist}
Let $n\ge2$, let $z_0,\dots,z_{n-1}$ be the vertices of a regular $n$-gon inscribed in the unit circle, and let $z$ be any point of that circle. Show that $\sum_{k=0}^{n-1}|z-z_k|^2=2n$.
\end{exercise}

\begin{exercise}[$\star\star$]\label{ex:cardano}
\emph{Cardano's formula, and the three cube roots.}
\begin{enumerate}[label=(\alph*)]
\item Let $p,q\in\R$. Show that if $u$ and $v$ satisfy $u^3+v^3=q$ and $uv=p/3$, then $x=u+v$ solves $x^3=px+q$. Deduce that $u^3$ and $v^3$ are the roots of $T^2-qT+p^3/27$, and hence Cardano's formula.
\item For $x^3=15x+4$, show that $u^3=2+11i$ and $v^3=2-11i$. Find all three cube roots of $2+11i$ in the form $(2+i)\omega^k$ with $\omega=e^{2i\pi/3}$. Pair each with the $v$ for which $uv=5$, and recover the three roots $4$, $-2+\sqrt3$ and $-2-\sqrt3$.
\end{enumerate}
\end{exercise}

\begin{exercise}[$\star\star$]\label{ex:equi}
\begin{enumerate}[label=(\alph*)]
\item Let $a,b,c$ be three distinct complex numbers with $|a|=|b|=|c|=1$ and $a+b+c=0$. Show that the triangle $abc$ is equilateral.
\item Let $abc$ be equilateral. Show that the triangle with vertices $a+b$, $b+c$ and $c+a$ is equilateral.
\end{enumerate}
\end{exercise}

\begin{exercise}[$\star\star$]\label{ex:cyclo}
\begin{enumerate}[label=(\alph*)]
\item Compute $\Phi_{10}$, $\Phi_{15}$ and $\Phi_{18}$ from Proposition~\ref{prop:cyclotomic}(a).
\item Show that $\Phi_{2n}(X)=\Phi_n(-X)$ for odd $n\ge3$.
\item Show that $\Phi_n(1)=p$ if $n$ is a power of a prime $p$, and that $\Phi_n(1)=1$ if $n$ has at least two distinct prime factors. Evaluate $(X^n-1)/(X-1)=\prod_{d\mid n,\,d>1}\Phi_d(X)$ at $X=1$, and use strong induction.
\end{enumerate}
\end{exercise}

\begin{exercise}[$\star\star\star$]\label{ex:napoleon}
\emph{Napoleon's theorem.} On the sides of a triangle $abc$, taken anticlockwise, build equilateral triangles pointing outward. Let $A$, $B$ and $C$ be their centres, with $A$ opposite $a$, so $A$ lies on the side $bc$, and cyclically for $B$ and $C$. Show that $ABC$ is equilateral.
\emph{Hint: the third vertex of the outward equilateral triangle on $bc$ is $b+(c-b)e^{-i\pi/3}$. Express $A$, $B$ and $C$, and apply Proposition~\ref{prop:equilateral}.}
\end{exercise}

\begin{exercise}[$\star\star\star$]\label{ex:cotsum}
\emph{A finite form of the Basel sum. Return to it after formal differentiation in Chapter~\ref{ch:polynomials}.} Let $n\ge2$, $\omega=e^{2i\pi/n}$ and $P(X)=1+X+\dots+X^{n-1}=\prod_{k=1}^{n-1}(X-\omega^k)$.
\begin{enumerate}[label=(\alph*)]
\item Show that $P'/P=\sum_{k=1}^{n-1}1/(X-\omega^k)$ and $(P'^2-P''P)/P^2=\sum_{k=1}^{n-1}1/(X-\omega^k)^2$, away from the roots, and evaluate both sums at $X=1$.
\item Show that $1/(1-\omega^k)=\tfrac12+(i/2)\cot(k\pi/n)$, and deduce
\[
\sum_{k=1}^{n-1}\cot^2\frac{k\pi}n=\frac{(n-1)(n-2)}3,\qquad\sum_{k=1}^{n-1}\frac1{\sin^2(k\pi/n)}=\frac{n^2-1}3.
\]
\item Deduce that $\sum_{k=1}^{n-1}1/|1-\omega^k|^2=(n^2-1)/12$. Using $\sin x\sim x$ as $x\to0$, guess what the identity says about $\sum_{k\ge1}1/k^2$ as $n\to\infty$. Chapter~\ref{ch:polynomials} turns the guess into a proof.
\end{enumerate}
\end{exercise}

\begin{challenge}
\begin{enumerate}[label=(\arabic*)]
\item \emph{The sign of the Gauss sum.} Compute $G=\sum_a\omega^{a^2}$ for $p=3,5,7,11,13$. The pattern is $G=\sqrt p$ when $p\equiv1\pmod4$, and $G=i\sqrt p$ when $p\equiv3\pmod4$. Gauss recorded the statement in his diary in the middle of May $1801$. The proof is the diary entry of $30$ August $1805$. Volume~5 gives a proof.
\item \emph{Constructible polygons.} The pentagon example expresses $\cos(2\pi/5)$ with one square root. The same attempt for $\cos(2\pi/7)$ leads to a cubic. Volume~3 explains the criterion: the regular $n$-gon is constructible with ruler and compass exactly when $n=2^a p_1\cdots p_r$ with distinct Fermat primes $p_i$. Gauss proved the sufficient direction; Wantzel, in $1837$, proved that the condition is necessary. Which $n\le30$ qualify?
\item \emph{Sums of roots of unity.} Find all triples of $n$th roots of unity $\alpha,\beta,\gamma$ with $\alpha+\beta=\gamma$. Draw the rhombus. For four unit complex numbers, a sum equal to $0$ splits into two antipodal pairs. The twelfth roots $e^{\pm i\pi/6}$ and $e^{\pm i5\pi/6}$ are one example: pair $e^{i\pi/6}$ with $e^{-i5\pi/6}$, and $e^{-i\pi/6}$ with $e^{i5\pi/6}$. The notes prove the three-term list and this splitting.
\end{enumerate}
Notes on the challenges are at the end of the chapter.
\end{challenge}

\begin{beauty}
The three cube roots of unity form an equilateral triangle, and the arrows cancel. That cancellation is the filter. If the exponent $k$ is not a multiple of $3$, the three numbers $P$ contributes at the cube roots of unity are a triangle of arrows, one for each copy of the coefficient $a_k$. The triangle closes, and $a_k$ drops out. What remains is a count. Exercise~\ref{ex:mod3} counts the subsets whose size is a multiple of $3$, and the cosine in the answer is the real part of a turn that survived. Put a square in the exponent instead of a linear exponent, and the same orthogonality proves $|G|^2=p$ in Exercise~\ref{ex:gauss}. The modulus is settled. The direction of $G$ is the part that took Gauss four years.
\end{beauty}

\section{Hints}

\begin{hint}{ex:lin}
Expand $\bigl((e^{i\theta}-e^{-i\theta})/(2i)\bigr)^4$ by the binomial theorem, and group conjugate terms into cosines. The double-angle formula is shorter.
\end{hint}
\begin{hint}{ex:binomcos}
The sum is the real part of $\sum_k\binom nk e^{ik\theta}=(1+e^{i\theta})^n$. Factor $1+e^{i\theta}=2\cos(\theta/2)\,e^{i\theta/2}$.
\end{hint}
\begin{hint}{ex:dirichlet}
The sum is geometric, with ratio $e^{i\theta}$. Multiply numerator and denominator by $e^{-i\theta/2}$.
\end{hint}
\begin{hint}{ex:mod3}
Apply Theorem~\ref{thm:filter} to $(1+X)^n$ with $m=3$ and $r=0$. Write $1+j$ and $1+j^2$ in exponential form.
\end{hint}
\begin{hint}{ex:sines}
Evaluate $(X^n-1)/(X-1)$ at $X=-1$. Account for the sign $(-1)^{n-1}$ when you pass from $-1-\omega^k$ to $1+\omega^k$. When $n$ is even, one factor is zero.
\end{hint}
\begin{hint}{ex:pentagon}
For (a), square $(\sqrt5-1)/4$. For (b), use $\sin^2+\cos^2=1$, and take the positive root: both angles lie in $(0,\pi/2)$.
\end{hint}
\begin{hint}{ex:circulant}
(a) The $r$th coordinate of $Cv_j$ is $\sum_s a_{(s-r)\bmod n}\,\omega^{sj}$. Substitute $t=s-r$ modulo $n$. (b) The vectors $v_j$ form a basis, because they are the columns of a Vandermonde matrix on distinct nodes. (c) Here $P=1+X$.
\end{hint}
\begin{hint}{ex:gauss}
$|G|^2=\sum_{a,b}\omega^{a^2-b^2}$. Substitute $a=b+h$. The exponent becomes $h^2+2bh$, which is linear in $b$. Sum in $b$ first, by Proposition~\ref{prop:orth}.
\end{hint}
\begin{hint}{ex:sqrtz}
Use $x^2-y^2=a$, $2xy=b$ and $x^2+y^2=|a+ib|$. For the quadratic, compute the discriminant and take its square root the same way.
\end{hint}
\begin{hint}{ex:sumdist}
$|z-z_k|^2=(z-z_k)(\bar z-\bar z_k)=2-2\operatorname{Re}(z\bar z_k)$, and $\sum_k\bar z_k=0$.
\end{hint}
\begin{hint}{ex:cardano}
(a) Expand $(u+v)^3$. (b) $|2+11i|=5\sqrt5=(\sqrt5)^3$, so every cube root has modulus $\sqrt5$. The partner is $v=5/u$. For $u=(2+i)\omega^k$ that partner is $(2-i)\omega^{-k}$.
\end{hint}
\begin{hint}{ex:equi}
(a) Expand $(a+b+c)^2$, and use $\bar a=1/a$ on the unit circle. (b) Apply criterion (ii) to the new vertices.
\end{hint}
\begin{hint}{ex:cyclo}
(a) Divide $X^n-1$ by the cyclotomic polynomials of the proper divisors. (b) If $\zeta$ is a primitive $n$th root and $n$ is odd, then $-\zeta$ has order $2n$. (c) The left side at $X=1$ equals $n$.
\end{hint}
\begin{hint}{ex:napoleon}
With $\alpha=e^{-i\pi/3}=-j$, the centre opposite $a$ is $A=\bigl((2-\alpha)b+(1+\alpha)c\bigr)/3$. Compute $A+jB+j^2C$, and use $1+j+j^2=0$.
\end{hint}
\begin{hint}{ex:cotsum}
(a) The logarithmic derivative of a product is the sum of the logarithmic derivatives. At $X=1$, one has $P(1)=n$, $P'(1)=n(n-1)/2$ and $P''(1)=n(n-1)(n-2)/3$. (b) Use $1-\omega^k=-2i\sin(k\pi/n)\,e^{ik\pi/n}$. The sum of the cotangents vanishes by the symmetry $k\mapsto n-k$.
\end{hint}

\section{Solutions}

\begin{solution}{ex:lin}
\idea Linearising means writing a power of sine as a sum of cosines of multiple angles. Those are easy to integrate. You can expand $\bigl((e^{i\theta}-e^{-i\theta})/(2i)\bigr)^4$. The double-angle formula is shorter, and it is the one written out here.
\writeup $\sin^2\theta=(1-\cos2\theta)/2$, so
\begin{align*}
\sin^4\theta&=\frac{1-2\cos2\theta+\cos^22\theta}4\eqbr=\frac{1-2\cos2\theta+\frac{1+\cos4\theta}2}4\eqbr=\frac{3-4\cos2\theta+\cos4\theta}8.
\end{align*}
For $k\ge1$, $\int_0^\pi\cos(2k\theta)\dd\theta=\bigl[\sin(2k\theta)/(2k)\bigr]_0^\pi=0$. Therefore $\int_0^\pi\sin^4\theta\dd\theta=3\pi/8$.
\end{solution}

\begin{solution}{ex:binomcos}
\idea The sum $\sum\binom nk\cos(k\theta)$ is the real part of the binomial expansion of $(1+e^{i\theta})^n$. The half-angle factorisation reads that real part off.
\writeup The sum is the real part of $\sum_k\binom nk e^{ik\theta}=(1+e^{i\theta})^n=\bigl(2\cos(\theta/2)\bigr)^n e^{in\theta/2}$. Hence
\[
\sum_{k=0}^n\binom nk\cos(k\theta)=2^n\cos^n\frac\theta2\,\cos\frac{n\theta}2.
\]
\end{solution}

\begin{solution}{ex:dirichlet}
\idea $D_n$ is a geometric sum with ratio $e^{i\theta}\ne1$. The target formula has half-angles in it, which is the signal to multiply by $e^{-i\theta/2}$.
\writeup Since $e^{i\theta}\ne1$,
\[
D_n(\theta)=e^{-in\theta}\,\frac{e^{i(2n+1)\theta}-1}{e^{i\theta}-1}.
\]
Multiply numerator and denominator by $e^{-i\theta/2}$:
\[
D_n(\theta)=\frac{e^{i(n+1/2)\theta}-e^{-i(n+1/2)\theta}}{e^{i\theta/2}-e^{-i\theta/2}}=\frac{\sin\bigl((n+\tfrac12)\theta\bigr)}{\sin(\theta/2)}.
\]
Pairing $k$ with $-k$ gives $D_n(\theta)=1+2\sum_{k=1}^n\cos(k\theta)$, so
\[
\sum_{k=1}^n\cos(k\theta)=\frac12\left(\frac{\sin\bigl((n+\tfrac12)\theta\bigr)}{\sin(\theta/2)}-1\right).
\]
Figure~\ref{fig:dirichlet} shows the graph: a peak of height $2n+1$ at $0$, and smaller oscillations elsewhere. Fourier series in Volume~5 are built on this peak.
\end{solution}

\begin{figure}[ht]
\centering
\fig{ch02-dirichlet}
\caption{The Dirichlet kernel $D_n(\theta)=\sin\bigl((n+\tfrac12)\theta\bigr)/\sin(\theta/2)$ for $n=3$ and $n=8$. The peak at $0$ has height $2n+1$, and the first zeros are at $\pm2\pi/(2n+1)$. The integral over $[-\pi,\pi]$ equals $2\pi$ for every $n$, because every $e^{ik\theta}$ with $k\ne0$ integrates to $0$ on that interval.}
\label{fig:dirichlet}
\end{figure}

\begin{solution}{ex:mod3}
\idea The count is $\sum_{3\mid k}\binom nk$. That is the filter of Theorem~\ref{thm:filter}, applied to $(1+X)^n$ with $m=3$ and $r=0$. The remaining work is to write $1+j$ as a single turn.
\writeup The number of such subsets is $\sum_{k\equiv0\pmod3}\binom nk$. Theorem~\ref{thm:filter} gives
\[
\frac13\bigl(2^n+(1+j)^n+(1+j^2)^n\bigr),
\]
where $j=e^{2i\pi/3}$. From $1+j+j^2=0$ we get $1+j=-j^2=e^{i\pi}e^{-2i\pi/3}=e^{i\pi/3}$, and $1+j^2=\overline{1+j}=e^{-i\pi/3}$. The count is
\[
\frac{2^n+2\cos(n\pi/3)}3.
\]
For $n=3$ this is $(8+2\cos\pi)/3=(8-2)/3=2$: the empty set, and the whole set.
\end{solution}

\begin{solution}{ex:sines}
\idea The product $\prod(1+\omega^k)$ is a value of $(X^n-1)/(X-1)$. At $X=-1$ that value is easy, and the even case is easier still, because $-1$ itself is an $n$th root of unity.
\writeup If $n$ is even, then $\omega^{n/2}=-1$, so one factor is $1+(-1)=0$ and the product is $0$.

If $n$ is odd, Proposition~\ref{prop:roots} gives $\prod_{k=1}^{n-1}(X-\omega^k)=(X^n-1)/(X-1)$. At $X=-1$,
\[
\prod_{k=1}^{n-1}(-1-\omega^k)=\frac{(-1)^n-1}{-2}=\frac{-2}{-2}=1.
\]
Then $1+\omega^k=(-1)(-1-\omega^k)$, so
\[
\prod_{k=1}^{n-1}(1+\omega^k)=(-1)^{n-1}\prod_{k=1}^{n-1}(-1-\omega^k)=1,
\]
because $n-1$ is even. For $n=3$ the factors are $1+j$ and $1+j^2$, and $(1+j)(1+j^2)=-j^2\cdot(-j)=j^3=1$.
\end{solution}

\begin{solution}{ex:pentagon}
\idea Part (a) is arithmetic with the value already proved. Part (b) is $\sin^2=1-\cos^2$, and the sign is the sign of sine on $(0,\pi/2)$.
\writeup (a) Let $c=(\sqrt5-1)/4$. Then
\[
2c^2-1=2\cdot\frac{5-2\sqrt5+1}{16}-1=\frac{3-\sqrt5}4-1=\frac{-1-\sqrt5}4.
\]
On the other side, $\cos(4\pi/5)=-\cos(\pi/5)=-({\sqrt5}+1)/4=(-1-\sqrt5)/4$. The two sides agree.

(b) Both angles lie in $(0,\pi/2)$, so both sines are positive.
\begin{align*}
\sin^2\frac\pi5&=1-\left(\frac{\sqrt5+1}4\right)^2=1-\frac{6+2\sqrt5}{16}=\frac{10-2\sqrt5}{16},\\
\sin\frac\pi5&=\frac{\sqrt{10-2\sqrt5}}4.
\end{align*}
The same computation with $(\sqrt5-1)/4$ gives $\sin(2\pi/5)=\sqrt{10+2\sqrt5}/4$. The side of the inscribed pentagon is $2\sin(\pi/5)=\sqrt{10-2\sqrt5}/2$.
\end{solution}

\begin{solution}{ex:circulant}
\idea The problem hands you the eigenvectors. Part (a) is a change of index, legal because the indices live modulo $n$ and $\omega^n=1$. For (b), $n$ eigenvectors that form a basis diagonalise $C$, and the determinant is the product of the eigenvalues. Part (c) is the case $P=1+X$.
\writeup (a) The $r$th coordinate of $Cv_j$ is, with $t=s-r$ taken modulo $n$,
\[
\sum_{s=0}^{n-1}a_{(s-r)\bmod n}\,\omega^{js}=\omega^{jr}\sum_{t=0}^{n-1}a_t\,\omega^{jt}=P(\omega^j)\,\omega^{jr}.
\]
The shift of index is legal because $\omega^{jn}=1$.

(b) The vectors $v_0,\dots,v_{n-1}$ are the columns of a Vandermonde matrix on the distinct numbers $1,\omega,\dots,\omega^{n-1}$. Its determinant is nonzero, by Proposition~\ref{prop:vdm} in Chapter~\ref{ch:linear}, so they form a basis. In this basis $C$ is diagonal with entries $P(\omega^j)$. The determinant does not depend on the basis, so it equals the product of those entries.

(c) Here $P=1+X$. From $X^n-1=\prod_j(X-\omega^j)$, the value at $X=-1$ is $\prod_j(-1-\omega^j)=(-1)^n-1$. Therefore
\begin{align*}
\det C&=\prod_{j=0}^{n-1}(1+\omega^j)=(-1)^n\bigl((-1)^n-1\bigr)\eqbr=1-(-1)^n,
\end{align*}
which is $2$ when $n$ is odd and $0$ when $n$ is even. When $n$ is even, the vector $(1,-1,1,\dots,-1)$ lies in the kernel, which is a second way to see that the determinant vanishes. Exercise~\ref{ex:sines} is the same product with the factor $1+\omega^0=2$ removed.
\end{solution}

\begin{solution}{ex:gauss}
\idea A square modulus is a double sum. The exponent $a^2-b^2$ is quadratic in both variables. The substitution $a=b+h$ makes it linear in $b$, and Proposition~\ref{prop:orth} kills every inner sum except one.
\writeup Since $\bar\omega=\omega^{-1}$,
\[
|G|^2=G\bar G=\sum_{a,b}\omega^{a^2-b^2},
\]
where $a$ and $b$ run through $\Z/p\Z$. The exponent only matters modulo $p$. Substitute $a=b+h$. For each fixed $h$ this is a bijection as $b$ runs through the field. Then $a^2-b^2=h^2+2hb$, and
\[
|G|^2=\sum_{h=0}^{p-1}\omega^{h^2}\sum_{b=0}^{p-1}\omega^{2hb}.
\]
By Proposition~\ref{prop:orth}, the inner sum equals $p$ when $p$ divides $2h$, and equals $0$ otherwise. Since $p$ is odd, $p\mid2h$ if and only if $p\mid h$, so only $h=0$ survives. Thus $|G|^2=\omega^0\cdot p=p$.
\end{solution}

\begin{solution}{ex:sqrtz}
\idea Square roots of $a+ib$ come from $x^2-y^2=a$, $2xy=b$ and $x^2+y^2=|a+ib|$. The quadratic formula then needs one square root of its discriminant, found the same way.
\writeup For $-7+24i$, the modulus is $25$, so $x^2=(25-7)/2=9$ and $y^2=(25+7)/2=16$. Since $2xy=24>0$, the two numbers have the same sign. The roots are $\pm(3+4i)$. For $2i$, one has $x^2=y^2=1$ and the same sign, so the roots are $\pm(1+i)$.

For $z^2-(1+i)z+i=0$, the discriminant is $(1+i)^2-4i=1+2i-1-4i=-2i$. The square roots of $-2i$ are $\pm(1-i)$, because $(1-i)^2=1-2i-1=-2i$. Therefore
\[
z=\frac{(1+i)\pm(1-i)}2,
\]
so $z=1$ or $z=i$. Check: $(z-1)(z-i)=z^2-(1+i)z+i$.
\end{solution}

\begin{solution}{ex:sumdist}
\idea Squared distances expand into a constant and a real part. The vertices of a regular polygon sum to $0$, by Proposition~\ref{prop:orth}, so the real parts disappear.
\writeup Since $|z|=|z_k|=1$,
\[
|z-z_k|^2=(z-z_k)(\bar z-\bar z_k)=2-2\operatorname{Re}(z\bar z_k).
\]
Sum over $k$. The vertices are $c\omega^k$ for some $c$ with $|c|=1$, and $\sum_k\omega^k=0$, so $\sum_k\bar z_k=0$. The sum of squared distances is $2n$.

For a check, take the regular hexagon and let $z$ be a vertex. The six squared distances are $0,1,3,4,3,1$, and they add to $12=2\cdot6$.
\end{solution}

\begin{solution}{ex:cardano}
\idea Part (a) is $(u+v)^3=u^3+v^3+3uv(u+v)$. Two numbers with a known sum and a known product are the roots of a quadratic. In (b) the formula produces $u^3=2+11i$, and there are three cube roots. Each pairs with exactly one $v$ through $uv=5$. Bombelli found the pair that gives $4$. The other two pairs give the other two real roots.
\writeup (a) Expand: $(u+v)^3=u^3+3u^2v+3uv^2+v^3=(u^3+v^3)+3uv(u+v)$. If $u^3+v^3=q$ and $3uv=p$, then $x=u+v$ satisfies $x^3=px+q$. The numbers $u^3$ and $v^3$ have sum $q$ and product $(uv)^3=p^3/27$, so they are the roots of $T^2-qT+p^3/27$, namely
\[
\frac q2\pm\sqrt{\frac{q^2}4-\frac{p^3}{27}}.
\]
Taking a cube root of each gives Cardano's formula. Each cube root has to be matched with the cube root of the other number that obeys $uv=p/3$. There are three legal pairs, and three roots of the cubic.

(b) With $p=15$ and $q=4$, one has $q^2/4-p^3/27=4-125=-121$, so $u^3=2+11i$ and $v^3=2-11i$. The computation in the text gives $(2+i)^3=2+11i$. The three cube roots are $u_k=(2+i)\omega^k$ for $k=0,1,2$, where $\omega=e^{2i\pi/3}=(-1+i\sqrt3)/2$. Indeed, if $u^3=u_0^3$ and $u_0\ne0$, then $(u/u_0)^3=1$.

The constraint $uv=p/3=5$ gives $v_k=5/u_k$. Since $5/(2+i)=2-i$, one has $v_k=(2-i)\omega^{-k}$. Hence
\[
x_k=u_k+v_k=(2+i)\omega^k+\overline{(2+i)\omega^k}=2\operatorname{Re}\bigl((2+i)\omega^k\bigr).
\]
For $k=0$, $x_0=4$. For $k=1$,
\[
(2+i)\omega=(2+i)\frac{-1+i\sqrt3}2=\frac{-2-\sqrt3+i(2\sqrt3-1)}2,
\]
so $x_1=-2-\sqrt3$. For $k=2$, $\omega^2=\bar\omega$, and $x_2=-2+\sqrt3$. The three roots add to $0$, which is the coefficient of $x^2$ in $x^3-15x-4$, as Vi\`ete's formulas require. Those formulas are Proposition~\ref{prop:viete} in Chapter~\ref{ch:polynomials}.
\end{solution}

\begin{solution}{ex:equi}
\idea Criterion (iii) asks for $a^2+b^2+c^2=ab+bc+ca$. The hypothesis $a+b+c=0$ turns that into a question about $ab+bc+ca$. On the unit circle, $\bar a=1/a$, and the sum $ab+bc+ca$ becomes a conjugate of $a+b+c$.
\writeup (a) Expand $0=(a+b+c)^2=a^2+b^2+c^2+2(ab+bc+ca)$. On the unit circle,
\[
ab+bc+ca=abc\Bigl(\frac1c+\frac1a+\frac1b\Bigr)=abc(\bar a+\bar b+\bar c)=abc\,\overline{a+b+c}=0.
\]
Thus $a^2+b^2+c^2=0=ab+bc+ca$. Proposition~\ref{prop:equilateral} says the triangle is equilateral. It is not degenerate, because the three points are distinct.

(b) Suppose $a+jb+j^2c=0$. The new vertices satisfy
\begin{align*}
(a+b)+j(b+c)+j^2(c+a)&=a(1+j^2)+b(1+j)+c(j+j^2)\\
&=-ja-j^2b-c=-j(a+jb+j^2c)=0,
\end{align*}
where $1+j^2=-j$, $1+j=-j^2$ and $j+j^2=-1$. The other orientation is the same computation with $j$ and $j^2$ exchanged. Geometrically, the new triangle is congruent to the old one: if $g=(a+b+c)/3$, the map $z\mapsto 3g-z$ sends $c,a,b$ to $a+b$, $b+c$, $c+a$.
\end{solution}

\begin{solution}{ex:cyclo}
\idea The factorisation $X^n-1=\prod_{d\mid n}\Phi_d$ is a recipe. The polynomial $\Phi_n$ is $X^n-1$ divided by the cyclotomic polynomials of the proper divisors. For (b) and (c), name the roots rather than expand.
\writeup (a) $\Phi_{10}=(X^{10}-1)/(\Phi_1\Phi_2\Phi_5)=(X^5+1)/(X+1)=X^4-X^3+X^2-X+1$. Next,
\[
\Phi_{15}=\frac{X^{15}-1}{\Phi_1\Phi_3\Phi_5}=\frac{X^{10}+X^5+1}{X^2+X+1}=X^8-X^7+X^5-X^4+X^3-X+1.
\]
The last step is the division carried out in the text for $\Phi_6$, one degree higher. Finally $\Phi_9=\Phi_3(X^3)=X^6+X^3+1$ by Proposition~\ref{prop:cyclotomic}(c), and part (b) gives $\Phi_{18}=\Phi_9(-X)=X^6-X^3+1$.

(b) Let $n$ be odd, and let $\zeta$ be a primitive $n$th root of unity. Then $(-\zeta)^{2n}=1$. If $(-\zeta)^m=1$, then $\zeta^{2m}=1$, so $n$ divides $2m$. Since $n$ is odd, $n$ divides $m$. Then $(-\zeta)^m=1$ forces $(-1)^m=1$, so $m$ is even, and in fact $2n$ divides $m$. Thus $-\zeta$ has order $2n$. The map $\zeta\mapsto-\zeta$ is injective. Also $\varphi(2n)=\varphi(n)$ for odd $n$: the integers in $\{1,\dots,2n\}$ that are coprime to $2n$ are the odd ones that are coprime to $n$, and there are $\varphi(n)$ of them. So the primitive $2n$th roots are exactly the numbers $-\zeta$. Hence
\[
\Phi_{2n}(X)=\prod(X+\zeta)=\prod\bigl(-(-X-\zeta)\bigr)=(-1)^{\varphi(n)}\Phi_n(-X).
\]
For $n\ge3$, $\varphi(n)$ is even, because the primitive $n$th roots come in distinct conjugate pairs $\zeta$ and $\zeta^{-1}$ (a root equal to its inverse would have order dividing $2$). The sign is $+1$, and $\Phi_{2n}(X)=\Phi_n(-X)$.

(c) At $X=1$, $(X^n-1)/(X-1)=n$, so $n=\prod_{d\mid n,\,d>1}\Phi_d(1)$. For $n=p$ this says $\Phi_p(1)=p$. For $n=p^k$, the divisors $d>1$ are $p,\dots,p^k$, and induction gives $p^k=p^{k-1}\Phi_{p^k}(1)$, so $\Phi_{p^k}(1)=p$. If $n$ has at least two distinct prime factors, induct on $n$. Every proper divisor $d>1$ that has two distinct prime factors contributes $\Phi_d(1)=1$. The prime-power divisors contribute $\prod_{p^j\mid n}p=\prod_{p\mid n}p^{v_p(n)}=n$, where $v_p(n)$ is the exponent of $p$ in $n$. Thus $n=n\cdot\Phi_n(1)$, and $\Phi_n(1)=1$. Checks: $\Phi_6(1)=1-1+1=1$, $\Phi_{10}(1)=1$, $\Phi_{12}(1)=1$.
\end{solution}

\begin{solution}{ex:napoleon}
\idea The centres are linear in $a$, $b$ and $c$, and the criterion of Proposition~\ref{prop:equilateral} is a linear condition. It is enough to watch the coefficient of each vertex and use $1+j+j^2=0$.
\writeup Let $\alpha=e^{-i\pi/3}$. Then $-j=e^{i\pi}e^{2i\pi/3}=e^{-i\pi/3}$, so $\alpha=-j$. For $abc$ anticlockwise, the third vertex of the outward equilateral triangle on $bc$ is $p=b+(c-b)\alpha$, and the centre is
\[
A=\frac{b+c+p}3=\frac{(2-\alpha)b+(1+\alpha)c}3=\frac{(2+j)b+(1-j)c}3.
\]
Cyclically, $B=\bigl((2+j)c+(1-j)a\bigr)/3$ and $C=\bigl((2+j)a+(1-j)b\bigr)/3$. The coefficient of $a$ in $3(A+jB+j^2C)$ is
\[
j(1-j)+j^2(2+j)=j-j^2+2j^2+1=1+j+j^2=0.
\]
The coefficient of $b$ is $(2+j)+j^2(1-j)=2+j+j^2-1=1+j+j^2=0$. The coefficient of $c$ is $(1-j)+j(2+j)=1-j+2j+j^2=1+j+j^2=0$. Thus $A+jB+j^2C=0$. Proposition~\ref{prop:equilateral} says that $ABC$ is equilateral and anticlockwise, for every triangle $abc$, including a degenerate one.

The theorem is traditionally named after Napoleon Bonaparte, who took an interest in mathematics. No document ties the theorem to him. The earliest printed problem we rely on is W.~Rutherford's, posed in \emph{The Ladies' Diary} for $1825$, with proofs in the $1826$ issue. Coxeter and Greitzer record that history.
\end{solution}

\begin{solution}{ex:cotsum}
\idea The logarithmic derivative of a product is the sum of the logarithmic derivatives of the factors. That is the product rule. Differentiating again produces the sum of squares. The half-angle formula turns $1/(1-\omega^k)$ into a cotangent.
\writeup (a) By the product rule, $P'=\sum_k\prod_{l\ne k}(X-\omega^l)$, so $P'/P=\sum_k 1/(X-\omega^k)$ away from the roots. Differentiate:
\[
\frac{P''P-P'^2}{P^2}=-\sum_k\frac1{(X-\omega^k)^2}.
\]
The second formula in the exercise is this identity multiplied by $-1$. At $X=1$, one has $P(1)=n$ and $P'(1)=\sum_{k=1}^{n-1}k=n(n-1)/2$. Also
\[
P''(1)=\sum_{k=2}^{n-1}k(k-1)=\frac{(n-1)n(n-2)}3,
\]
by writing $k(k-1)=2\binom k2$ and summing, or by telescoping as in Chapter~\ref{ch:sums}. Therefore
\begin{align*}
\sum_{k=1}^{n-1}\frac1{1-\omega^k}&=\frac{n-1}2,\\
\sum_{k=1}^{n-1}\frac1{(1-\omega^k)^2}&=\frac{(n-1)^2}4-\frac{(n-1)(n-2)}3=\frac{(n-1)(5-n)}{12}.
\end{align*}

(b) Let $\theta_k=k\pi/n$. Then $1-\omega^k=-2i\sin\theta_k\,e^{i\theta_k}$, so
\[
\frac1{1-\omega^k}=\frac{e^{-i\theta_k}}{-2i\sin\theta_k}=\frac12+\frac i2\cot\theta_k.
\]
The real part $1/2$ matches the first sum in (a). Square:
\[
\frac1{(1-\omega^k)^2}=\frac14-\frac14\cot^2\theta_k+\frac i2\cot\theta_k.
\]
Sum over $k$. The imaginary parts cancel in pairs, because $\cot\theta_{n-k}=-\cot\theta_k$. Compare with (a):
\[
\frac{n-1}4-\frac14\sum_k\cot^2\theta_k=\frac{(n-1)(5-n)}{12},
\]
so $\sum_{k=1}^{n-1}\cot^2(k\pi/n)=(n-1)(n-2)/3$. Add $n-1$, using $1/\sin^2=1+\cot^2$:
\[
\sum_{k=1}^{n-1}\frac1{\sin^2\theta_k}=\frac{(n-1)(n-2)}3+(n-1)=\frac{n^2-1}3.
\]

(c) Since $|1-\omega^k|^2=4\sin^2\theta_k$,
\[
\sum_{k=1}^{n-1}\frac1{|1-\omega^k|^2}=\frac14\cdot\frac{n^2-1}3=\frac{n^2-1}{12}.
\]
For $n=7$ the sum is $48/12=4$. Divide this identity by $n^2$:
\[
\frac1{n^2}\sum_{k=1}^{n-1}\frac1{|1-\omega^k|^2}=\frac{n^2-1}{12n^2},
\]
which tends to $1/12$. The factor $4$ in $|1-\omega^k|^2=4\sin^2(k\pi/n)$ stays in the approximation. For fixed $k$, $\sin(k\pi/n)\sim k\pi/n$, so
\[
\frac1{n^2|1-\omega^k|^2}=\frac1{4\,n^2\sin^2(k\pi/n)}\sim\frac1{4\pi^2 k^2}.
\]
The terms with $k$ near $n$ mirror the small ones. Reading the limit $1/12$ off those pairs gives
\[
2\sum_{k\ge1}\frac1{4\pi^2 k^2}=\frac1{12},
\]
and therefore $\sum_{k\ge1}1/k^2=\pi^2/6$. Passing the limit inside the sum is the content of Exercise~\ref{ex:basel}.
\end{solution}

\subsection*{Notes on the challenges}

(1) The five values are $G=i\sqrt3$, $\sqrt5$, $i\sqrt7$, $i\sqrt{11}$ and $\sqrt{13}$, up to the usual floating-point dust if you compute them numerically. Gauss's diary entry $118$, from the middle of May $1801$, records the sign. Entry $123$, dated $30$ August $1805$, records the proof. The phrase about a flash of lightning is not in the diary. It is in his letter to Olbers of $3$ September $1805$. The published proof is the \emph{Summatio quarumdam serierum singularium} of $1811$. Volume~5 proves the sign. This exercise proves only the modulus.

(2) The admissible $n\le30$ are $1$, $2$, $3$, $4$, $5$, $6$, $8$, $10$, $12$, $15$, $16$, $17$, $20$, $24$ and $30$. Equivalently, $\varphi(n)$ is a power of $2$, which happens when $n=2^a p_1\cdots p_r$ with distinct Fermat primes $p_i$. The five known Fermat primes are $3$, $5$, $17$, $257$ and $65537$. Whether another exists is open. For $n=7$, one has $\varphi(7)=6$. Let $\mu=2\cos(2\pi/7)$. The minimal polynomial of $\mu$ is $T^3+T^2-2T-1$. The only possible rational roots are $\pm1$, and neither is a root, so the cubic is irreducible over $\Q$. Thus $\cos(2\pi/7)=\mu/2$ has degree $3$, which is not a power of $2$. The same number is a root of $8x^3+4x^2-4x-1=0$. Volume~3 proves that a constructible number has degree a power of $2$, and proves Gauss's sufficient condition as well.

(3) Suppose $\alpha+\beta=\gamma$ with $|\alpha|=|\beta|=|\gamma|=1$. Divide by $\gamma$. The new equation is $u+v=1$ with $|u|=|v|=1$. Then $|1-u|^2=1$. Expand: $2-2\operatorname{Re}u=1$, so $\operatorname{Re}u=1/2$. A unit complex number with real part $1/2$ is $e^{\pm i\pi/3}$, and the two choices are partners, because $e^{i\pi/3}+e^{-i\pi/3}=1$. So $\{\alpha,\beta\}=\{\gamma e^{i\pi/3},\gamma e^{-i\pi/3}\}$. These three numbers are $n$th roots of unity exactly when $e^{i\pi/3}$ is an $n$th root of unity, that is, when $6$ divides $n$. There is no second family. If $u=v$, then $2u=1$ and $|u|=1/2$, which is off the circle.

Four unit complex numbers that sum to $0$ split into two antipodal pairs. Rotate so that one of them is $1$, and write the next two as $e^{i\alpha}$ and $e^{i\beta}$. The fourth has modulus $1$, so $|1+e^{i\alpha}+e^{i\beta}|=1$. Expand the square:
\[
|1+e^{i\alpha}+e^{i\beta}|^2=3+2\bigl(\cos\alpha+\cos\beta+\cos(\alpha-\beta)\bigr).
\]
The left side equals $1$, so
\[
\cos\alpha+\cos\beta+\cos(\alpha-\beta)=-1.
\]
Set $u=(\alpha+\beta)/2$ and $v=(\alpha-\beta)/2$. Then $\cos\alpha+\cos\beta=2\cos u\cos v$ and $\cos(\alpha-\beta)=2\cos^2 v-1$, so the cosine equation becomes $2\cos v\,(\cos u+\cos v)=0$. Thus $\cos v=0$ or $\cos u=-\cos v$.

If $\cos v=0$, then $\alpha-\beta$ is an odd multiple of $\pi$, and $e^{i\alpha}=-e^{i\beta}$. Those two sum to $0$, so the remaining number is $-1$.

If $\cos u=-\cos v$, then $u=(2k+1)\pi\pm v$ for some integer $k$. With the minus sign, $\alpha=(2k+1)\pi$. With the plus sign, $\beta=(2k+1)\pi$. Either way one of those two numbers is $-1$, and the other two sum to $0$.

So the four numbers are $\alpha_0$, $-\alpha_0$, $\gamma$ and $-\gamma$. The twelfth roots $e^{\pm i\pi/6}$ and $e^{\pm i5\pi/6}$ are one such splitting. Pair $e^{i\pi/6}$ with $e^{-i5\pi/6}$, and $e^{-i\pi/6}$ with $e^{i5\pi/6}$: each second number is the first multiplied by $-1$. Mann, in $1965$, classified minimal linear relations among roots of unity, relations in which no nonempty proper subsum vanishes. For four unit numbers, the splitting above is the whole list.

\section*{Sources and further reading}

\begin{itemize}\raggedright
\item T.~Needham, \emph{Visual Complex Analysis}, 25th anniversary ed., Oxford University Press, 2023, Chapters~1--2. Complex numbers as geometry.
\item T.~Andreescu and D.~Andrica, \emph{Complex Numbers from A to\dots Z}, 2nd ed., Birkh\"auser, 2014. Roots of unity, and problems in the style of Exercises~\ref{ex:pentagon}, \ref{ex:equi} and~\ref{ex:napoleon}.
\item R.~L.~Graham, D.~E.~Knuth and O.~Patashnik, \emph{Concrete Mathematics}, 2nd ed., Addison-Wesley, 1994, Chapter~7. The roots-of-unity filter.
\item T.~H.~Cormen, C.~E.~Leiserson, R.~L.~Rivest and C.~Stein, \emph{Introduction to Algorithms}, 4th ed., MIT Press, 2022, Chapter~30. The fast Fourier transform, which is the filter made fast. Volume~5 of this series proves it.
\item B.~C.~Berndt, R.~J.~Evans and K.~S.~Williams, \emph{Gauss and Jacobi Sums}, Wiley, 1998. The sum $G$ of Exercise~\ref{ex:gauss}, including the sign.
\item P.~J.~Nahin, \emph{An Imaginary Tale: The Story of $\sqrt{-1}$}, Princeton University Press, 1998. Cardano, Bombelli, Wessel, Argand and Gauss, with the sources. Nahin records Wessel's paper as presented in $1797$, printed in $1799$, noticed again in $1895$, and translated into French in $1897$.
\item G.~Cardano, \emph{Ars Magna, or The Rules of Algebra}, translated by T.~R.~Witmer, MIT Press, 1968; Dover reprint $1993$. Chapter~37 has the sum $10$ and the product $40$. Witmer renders \emph{dimissis incruciationibus} as putting aside the mental tortures, and a footnote records the other reading.
\item R.~Bombelli, \emph{L'Algebra}, Bologna, $1572$. The phrase ``pi\`u tosto sofistica che reale'' and the treatment of $x^3=15x+4$ are discussed in J.~Stillwell, \emph{Mathematics and Its History}, 3rd ed., Springer, $2010$, Chapter~14, and in Nahin. Stillwell also records Brahmagupta's composition of sums of two squares.
\item F.~Cajori, \emph{A History of Mathematical Notations}, vol.~2, Open Court, $1929$, \S498. Euler's symbol $i$, in a memoir presented on $5$ May $1777$ and printed in $1794$.
\item J.~J.~Gray, \emph{A commentary on Gauss's mathematical diary, 1796--1814}, Expositiones Mathematicae~2 (1984), 97--130. Entry~1, $30$ March $1796$, is the $17$-gon. The entries of May $1801$ and $30$ August $1805$ are the Gauss sum. Gray points to the letter to Olbers of $3$ September $1805$ for Gauss's account of the proof.
\item E.~Lehmer, \emph{On the magnitude of the coefficients of the cyclotomic polynomial}, Bulletin of the American Mathematical Society~42 (1936), 389--392. Schur's unpublished argument of $1931$ that the coefficients are unbounded, and Migotti's coefficient $-2$ in $\Phi_{105}$.
\item H.~S.~M.~Coxeter and S.~L.~Greitzer, \emph{Geometry Revisited}, Mathematical Association of America, $1967$, Section~3.3. Napoleon's theorem, and Rutherford's problem of $1825$.
\item H.~B.~Mann, \emph{On linear relations between roots of unity}, Mathematika~12 (1965), 107--117. Minimal vanishing sums of roots of unity.
\end{itemize}
